% Section 8, threats to validity (restructure plan, part 2): ordered by their effect on the lead result, each tagged by
% the runs it concerns. Every number is a macro (\fv, \plv, \dv; their ledgers). The pilot detail of the reference-check
% review and of contract recovery stays with the construction pilot's material. The item on optimisation pressure was
% added in a later round (decision-log entry D82), after a critical read; the fresh attempts' inputs are as
% Section sec:design-tasks gives them (protocol §15.4), the held-out discipline's exceptions are App. B's
% (app:deviations), and a clause of the misuse paragraph of Section sec:ethics points here. Changed in a later round
% (decision-log entry D83), after reviews by model agents: "the zeros of C3, C4 and B2" became their acceptance of no
% violator, with its macros (cv.c3, cv.c4, cv.b2), since their false acceptance is c3.fa, not zero; and the labels item
% says that both label sources were validated on the reference.
\section{Threats to validity}
\label{sec:threats}
\begin{enumerate}
\item \textbf{Labels} [\texttt{\FRESHRUN}, pilot, natural]. No human label exists; every rate rests on the label rule's model proposals, which differ by candidate source (Table~\ref{tab:label-rules}): the ambiguity rule applies to constructed candidates, and an attempt is invalid if a reference check or the regression norm fails. In CS-02, \fv{unspec.n} attempts pass every reference check and fail a base test: they are invalid under the label rule, the second labeler calls them valid, and they are the whole false acceptance of C3, C4 and B2. A violation the reference checks miss is ambiguous for a constructed pilot candidate and valid for an attempt, so C3's, C4's and B2's acceptance of no violator (\fv{cv.c3}, \fv{cv.c4}, \fv{cv.b2}) holds only for what those checks exercise (\S\ref{sec:pilot}). The reference checks and the correction-aware verifiers share a model family, the reference diff and requirement lists that cover the same requirements, so C3's and C4's acceptance of no violator is partly agreement between renderings of the same requirements (\S\ref{sec:design-labels}). The second labeler's agreement is agreement between models of the same model family that read the same reference contract and whose checks were both validated on the reference.
\item \textbf{Asymmetric reference-check review} [\texttt{\FRESHRUN}, pilot]. Only CS-02's reference checks went through the review-and-fix loop (D33); the others were frozen, and used for labels, before that loop existed. The checks that label the silent-loss attempts of CS-03 and CS-04 were never reviewed.
\item \textbf{One model family of one provider; scenarios by the study's agent} [all]. By the study's procedure, every model-written artifact, the labels included, comes from models of one model family of one provider, which can correlate errors across conditions and with the labels; for several roles no record holds the identifier, so the release cannot confirm it (Table~\ref{tab:roles}). The \fv{families.w} controlled scenarios were written by the study's agent and registered within \dv{design.scenarios.span} minutes on \dv{design.scenarios.date} (from the CS-01 registration that superseded its first; D5), not drawn from a sampling frame.
\item \textbf{Few families} [\texttt{\FRESHRUN}, pilot]. No family-level sign-flip test can reach a one-sided $p$ below \fv{cmp.floor.one}, so there is no power for any effect. The pilot's planned bootstrap intervals describe variability; for a rate of \plv{c3.fa}, the upper end of the exact \dv{design.exact.level}\% binomial interval is \plv{c3.fa.exact.upper}.
\item \textbf{The fresh stratum is not a new sample of tasks} [\texttt{\FRESHRUN}]. It reuses the pilot's tasks and frozen verifiers and was designed after the pilot's results were known; $k=\dv{design.k}$ attempts per configuration and variant say little about variance, and no further run repeats them. The attempts' context held harness material (\S\ref{sec:design-tasks}), and an attempt may infer that it is part of a study.
\item \textbf{The information lost depends on the contract proposer} [pilot, \texttt{\FRESHRUN}]. For CS-02 the initial-only contract already holds, as inferred requirements, requirements covering everything the correction later stated (Table~\ref{tab:coverage}), as the initial-only contracts of NP-05 and NP-18 cover their correction requirements (Table~\ref{tab:nr-coverage}). The initial-only arm measures what the study's proposer recovers from the first request, not a fixed property of the session.
\item \textbf{Controlled population} [\texttt{\FRESHRUN}, pilot]. The study wrote the user turns and designed each session to contain a correction, which reached the session agent with a short harness frame; this feedback is synthetic.
\item \textbf{No optimisation pressure} [all]. No policy was optimised against any of the study's verifiers: every candidate that the rates count is an honest attempt at the task (a fresh attempt, which receives no verifier, \S\ref{sec:design-tasks}; a session's pre-correction state; a pull request's pre-review state) or a constructed alternative or near miss, built under the held-out discipline whose exceptions Appendix~\ref{app:deviations} lists. The false-acceptance rates, C3's, C4's and B2's acceptance of no violator (\fv{cv.c3}, \fv{cv.c4}, \fv{cv.b2}) included, therefore say nothing about a verifier used as a reward under optimisation, where a policy can search out what the verifier does not check.
\item \textbf{Natural arm} [natural]. S9's verdicts are model-judged screening yields, not prevalence; a reviewer's comment on a pull request is not an in-session correction from the requester; the reference is not an oracle; the run is below the size rule, an automated pilot, reported as case studies (\S\ref{sec:natural-cases}). Some of the natural arm's builder decisions await the owner's review, and the others took effect with no recorded review by a person; Appendix~\ref{app:deviations} lists them.
\item \textbf{B1's exploit surface} [all]. B1 imports candidate code into the reporting process, so a deliberately adversarial candidate can forge its report; B1 is meaningful only for non-adversarial candidates (\S\ref{sec:repro}).
\item \textbf{Contamination} [all]. The source repositories and the released tasks are public and may enter training data, and public merged fixes may be memorised. Protocol~\S15.6 states that the natural arm's pull requests and their repositories predate the generators' training data; with no model named, the paper cannot check this.
\item \textbf{The literature} [paper]. For patch overfitting in program repair, the test oracle problem and requirements change, our related work cites entry points only (\S\ref{sec:related}); it does not survey those literatures.
\end{enumerate}
